% Draft subsection for the Methods section of main.tex.
% Suggested placement: after \subsection{Average-face generation},
% before \subsection{Landmark visualization}.
%
% TODO before merging into main.tex:
%   - replace citation keys ddcolor / gfpgan / stylegan3 with the real .bib keys
%   - optionally add the per-cohort image-count range (approx. 60-400)
%   - the landmark-based crop/resize function is written but not yet in the
%     repository; add the command name once it is merged
%   - figure image/synthetic_image_generation.png is not referenced yet

\subsection{Synthetic Facial Image Generation}

FaceKit additionally supports the generation of synthetic facial images.
Two usage modes are provided. Users with a sufficiently large image
collection can train a generator on their own cohort and sample from the
resulting weights, while users without access to such data can sample
directly from generators pre-trained on ten rare-disease cohorts and
distributed with FaceKit.

Facial photographs collected for rare-disease research are often
retrospective and heterogeneous in acquisition quality, and the images used
for training are therefore standardized beforehand. Images that cannot be
used for morphological modelling may optionally be excluded at this stage,
including those obscured by privacy blocks, those of extremely low quality,
and those in which facial features are not clearly visible. Grayscale or
colour-degraded photographs are then colorized using DDColor~\cite{ddcolor},
and blurred or low-resolution faces are restored and upsampled using
GFPGAN~\cite{gfpgan}. Finally, FaceKit locates the facial region from the
detected landmarks, removes the surrounding background, and resizes the
retained facial area to $256\times256$ pixels, so that all training images
share a common facial framing and resolution. This preparation applies to
the training corpus only and is not required when sampling from an existing
generator.

Generators are trained using StyleGAN3~\cite{stylegan3}. A latent vector
$z \sim \mathcal{N}(0, I)$ is transformed by a mapping network into an
intermediate style vector, which modulates the convolutions of a synthesis
network to produce an image, while a discriminator is trained in parallel to
distinguish real from generated images; the two networks are optimized
adversarially. Each cohort, comprising on average more than one hundred
prepared images, is trained separately at $256\times256$ resolution with an
R1 regularization weight of $\gamma = 2$, horizontal-flip augmentation
enabled to increase effective data diversity, and a training length of
5{,}000 kimg. The generator and discriminator learning rates are initialized
at $2.5\times10^{-3}$ and $2\times10^{-3}$, respectively, and follow a
cosine decay schedule. Synthetic images are then obtained by drawing latent
vectors from the standard normal distribution and passing them through
either a user-trained or a pre-trained generator.
